% This chapter should orient a reader who is knowledgeable in the broad
% discipline but may not be a specialist in the precise dissertation topic.
%
% Depending on the field and dissertation structure, it may:
% - introduce the broad research context;
% - explain why the problem matters;
% - define important terminology;
% - synthesize the most relevant prior work;
% - identify unresolved questions or limitations;
% - state the dissertation's research questions, hypotheses, or objectives;
% - explain the overall methodological approach;
% - describe how the chapters fit together.
%
% A separate literature-review or methods chapter may make some of these
% elements unnecessary here.

\chapteropen
\chapter{Introduction}\label{ch:introduction}
\quoteOne




In this chapter, we outline the motivation and scientific scope of this thesis, summarizing the key questions addressed, before introducing the context required to frame the analysis. 

\section{Research context}

These section titles are illustrative and can (and should!) be changed.

Nevertheless, an introduction should generally do more than present background. By the end, the reader should understand

\begin{itemize}
    \item what problem is being addressed;
    \item why it is worth addressing;
    \item what the dissertation contributes;
    \item how the rest of the work is organized.
\end{itemize}

\subsection{Subsections (level 3) can also be used}

\section{Motivation and open questions}

\section{Dissertation aims}

\section{Organization of the dissertation}


Following this introduction, this template is structured in the following way. An example chapter is included in \cref{ch:chapter2}, where the layout is closer to a conventional research chapter. In \cref{ch:paper-example}, we analyze the distribution of \glspl{gc} around major galaxies in the Coma Cluster, focusing on systems with anomalously low \gls{gc} population counts. We then quantify the azimuthal asymmetry of these deficient \glspl{gcs} to search for signatures of environmental stripping. In \cref{ch:chapter4}, we extend this analysis by using \gls{gc} color as a proxy for metallicity, testing whether the blue and red \gls{gc} subpopulations respond differently to environmental processing in the dense cluster core. In \cref{ch:discussion}, we synthesize the results from these three studies and discuss their implications for the evolution of galaxies in dense cluster environments. Finally, in \cref{ch:conclusions}, we summarize the main findings of the thesis and place them in the broader context of \glspl{css}, environmental processing, and galaxy evolution.

